\documentclass[10pt,a4paper]{amsart}
\usepackage[margin=19mm]{geometry}
\usepackage{amsmath,amssymb,mathtools}
\usepackage[scr=rsfs,cal=euler]{mathalfa}
\usepackage{xfrac}
\usepackage[unicode,hidelinks,bookmarks=false]{hyperref}
\newcommand{\ie}{i.e.\ }
\newcommand{\cf}{cf.\ }
\newcommand{\resp}{respectively\ }
\newcommand{\Subsection}{\subsection}
\providecommand{\nobreakdash}{\nobreak\hbox{-}}
\setlength{\emergencystretch}{2em}
\multlinegap0pt
\usepackage{amssymb}
\usepackage{dsfont}
\usepackage{xcolor}
\newtheorem{lemma}{Lemma}
\title{Delta Theory of Anderson Modules I: Differential Characters}
\author{Sudip Pandit, Arnab Saha}
\date{Source: arXiv:2212.02137v2 (CC BY 4.0)}
\begin{document}
\maketitle

\noindent\emph{Source and licence.} Delta Theory of Anderson Modules I: Differential Characters. Version arXiv:2212.02137v2 (CC BY 4.0), distributed by arXiv under the Creative Commons Attribution 4.0 International licence, http://creativecommons.org/licenses/by/4.0/, as recorded in the arXiv licence metadata for this exact version and shown on the abstract page https://arxiv.org/abs/2212.02137v2. Full licence notice: \texttt{licenses/arxiv-2212.02137-CC-BY-4.0.txt}.\par\smallskip

\noindent\textit{Editorial scope.} Counted unit: the article's lemma esti (source0.tex lines 1093--1094). The article's complete original proof of it is delivered with it (source0.tex lines 1095--1126, one \texttt{proof} environment of 32 lines). No mathematics is written, repaired or re-derived: the statement and the proof are the article's own lines, in the article's own notation, and no retained line is substituted or re-flowed.  No further source-local context is needed: every cross-reference of every retained line resolves inside the two delivered spans.  Nothing was omitted from any retained span, so the package carries no omission record.  No retained line contains a citation, so the wrapper carries no bibliography and no bibliography entry.  No retained line contains a figure environment, an image-inclusion command or drawing code, so no image is delivered and the package declares no asset dependency.  Editorial formatting only, in addition: an AMS \texttt{amsart} preamble that restates the article's own declarations and macros used by the retained lines, namely the environments \texttt{lemma} on separate counters, together with the standard packages those lines need.  The retained environments print in this document as Lemma 1 in order of appearance, whereas the article prints the same items with the numbering of its own full document.  The complete native source of the article as distributed in the arXiv e-print for this exact version is bundled unchanged as \texttt{sources/134468/source0.tex}.
\begin{lemma}\label{esti} Let $r_{0}$ be the smallest number such that $q^{r_{0}}>d$. Let $L_{0}\geq r_{0}d$ be an integer. Suppose for each positive integer $k$,  there is an integer  $L_{k}$ satisfying $$L_{k}\geq \min~\{L_{k-i}+(q^{i}-1)q^{k-i}-d, ~\mathrm{for}~i=1,\ldots k\}.$$ Then $L_{k}\geq 0$ and $L_{k}\rightarrow \infty$ as $k\rightarrow \infty$.   
\end{lemma}\color{black}

\begin{proof}
We will first prove that for each $k$, $L_{k}\geq 0$.

\noindent \textbf{\underline{Case-I}:} Let $k\leq r_{0}$. We claim that $L_{k}\geq (r_{0}-k)d$, and we will show it by induction. By our assumption $L_{0}\geq (r_{0}-0)d=r_{0}d$. Assume that the statement holds for numbers $j<k$. Then for $i=1,\ldots k$, we have 

\begin{align*}
L_{k-i}+(q^{i}-1)q^{k-i}-d &\geq L_{k-i}-d\\
&\geq (r_{0}-(k-i))d-d\\
&\geq (r_{0}-k+i-1)d\\
&\geq (r_{0}-k)d\\
&\geq 0.
\end{align*}
  
\noindent\textbf{\underline{Case-II}:} Let $k> r_{0}$. Then by our choice of $r_{0}$, we have $q^{r_{0}}>d$. Hence for $i=1,\ldots k$, we have 
\begin{align*}
L_{k-i}+(q^{i}-1)q^{k-i}-d&\geq (q^{k}-q^{k-i})-d\\
&\geq (q^{k}-q^{k-1})-d\\
&= q^{k-1}(q-1)-d\\
&\geq q^{r_{0}}-d\\
&\geq  0.
\end{align*}
Note that for $k>r_{0}$, the above estimate gives us
\begin{align*}
L_{k-i}+(q^{i}-1)q^{k-i}-d\geq q^{k-1}(q-1)-d,~ \mathrm{for~ all}~ i=1,\ldots k.
\end{align*}
Hence 
\begin{align*}
L_{k}&=\min \{L_{k-i}+(q^{i}-1)q^{k-i}-d ~|~ i=1,\ldots k\}\\
&\geq q^{k-1}(q-1)-d.
\end{align*}
Therefore $L_{k}\rightarrow \infty$ as $k\rightarrow \infty$ and that completes the proof.
\end{proof}

\end{document}
